
\section{Data Generation}
Algorithm~\ref{algo:data_generation} outlines the automatic pipeline for generating high-quality multi-agent collaborative reasoning data used in M500. Starting from a raw pool of questions, the procedure filters for interdisciplinary and sufficiently complex problems using an LLM. For each qualified question, the MAS generates a reasoning trace. The resulting trace is included in the dataset only if it satisfies predefined quality criteria, including consensus, correct formatting, and correctness. This process continues until the target dataset size is reached.

\begin{figure}[h]
\begin{minipage}{\linewidth}
\begin{algorithm}[H]
\small
\caption{MAS Collaborative Data Generation}
\label{algo:data_generation}
\begin{algorithmic}[1]
\State \textbf{Input:} Question pool $\mathcal{Q}_{\text{raw}}$, LLM $f$, Multi-Agent System $\text{MAS}$, Target size $N$
\State \textbf{Output:} High-quality dataset $\mathcal{D}$

\Procedure{GenerateData}{}
    \State $\mathcal{D} \gets \emptyset$
    \ForAll{$q \in \mathcal{Q}_{\text{raw}}$}
        \If{not \Call{IsInterdisciplinary}{$f, q$} \textbf{or} \Call{TokenUsage}{$f(q)$} $<$ 1024}
            \State \textbf{continue}
        \EndIf
        \State $t \gets$ \Call{MASReasoningTrace}{$\text{MAS}, f, q$}
        \If{\Call{IsValidTrace}{$t$}}
            \State $\mathcal{D} \gets \mathcal{D} \cup \{(q, t)\}$
        \EndIf
        \If{$|\mathcal{D}| = N$}
            \State \textbf{break}
        \EndIf
    \EndFor
    \State \Return $\mathcal{D}$
\EndProcedure
\\
\Function{IsInterdisciplinary}{$f, q$}
    \State \Return $f$ predicts $q$ requires multiple experts to accomplish
\EndFunction
\\
\Function{TokenUsage}{$f(q)$}
    \State \Return Number of tokens used in $f$'s answer to $q$
\EndFunction
\\
\Function{IsValidTrace}{$t$}
    \State \Return $t$ satisfies consensus, format compliance, and correctness
\EndFunction
\end{algorithmic}
\end{algorithm}
\end{minipage}
\end{figure}

\section{Prompts}\label{sec:appendix_prompts}
To support clarity, we provide the full set of prompts used by each agent in the AgentVerse framework on mathematical reasoning tasks. Each agent role—\textit{CEO}, \textit{Expert Recruiter}, \textit{Problem Solver 1}, \textit{Problem Solver 2}, and \textit{Evaluator}—is governed by a system prompt and a user prompt that define its responsibilities, behavior, and expected outputs. The prompts are carefully designed to simulate realistic collaboration and maintain strict adherence to role-specific constraints.





\begin{tcolorbox}[
  title=CEO Prompts,
  breakable,
  colback=white,
  colframe=black
]

\textbf{System Prompt} \\
You are the CEO of a collaborative problem-solving system. Your responsibilities include: \\
1. Monitoring solution progress and resource allocation \\
2. Making strategic decisions about continuation/termination \\
3. Managing expert recruitment and retention \\
4. Directing discussion focus areas when the solution is not correct \\
5. Adjusting reasoning depth through token budgets \\

Previous system state: \\
- Task: \$\{task\_description\} \\
- Latest solution: \$\{current\_solution\} \\
- Evaluation feedback: \$\{evaluation\_feedback\} \\
- Current resources: \$\{current\_resources\} \\

\textbf{User Prompt} \\
Now, you need to decide the system state for this round. Carefully consider the following: \\
- Choose \textless Stop\textgreater\ only if solution is correct \\
- Recruit experts based on skill gaps identified in evaluation and do not recruit more than 4 experts, typically only 2-3 agents are needed for ordinary tasks and 4 agents are needed for complex tasks \\
- Direct discussion to address weakest solution aspects \\
- Set token budget proportional to the task complexity, token usages should choose from [0, 2048, 4096, 8192, 16384, 32000], typically 2048 tokens for simple tasks, 8192 tokens for tasks require medium reasoning, and 16384 or more tokens for complex reasoning tasks \\

Your response must strictly follow this structure: \\
\#\#\# Decision: \textless Continue\textgreater\ or \textless Stop\textgreater\ \\
\#\#\# Recruit Number: Number of experts to recruit in this round, should be an integer between 1 and 4 \\
\#\#\# Direction: Discussion direction based on the task description, latest solution, critic opinions, and evaluation feedback \\
\#\#\# Maximum Tokens: Maximum tokens for each agent in this round, should be an integer between 2048 and 32000 \\

\end{tcolorbox}





\begin{tcolorbox}[
  title=Expert Recruiter Prompts,
  breakable,
  colback=white,
  colframe=black
]

\textbf{System Prompt} \\
\# Role Description \\
You are the leader of a group of experts, now you are facing a math problem: \\
\$\{task\_description\} \\

\# Primary Objective \\
Your sole responsibility is to recruit \$\{cnt\_critic\_agents\} experts in different specialized fields to solve the math problem. \\
- DO NOT attempt to solve the problem yourself \\
- DO NOT propose any solutions or calculations \\

\# Recruitment Focus \\
Your selection should be based on: \\
1. Identifying which expertise domains are relevant to this math problem type \\
2. Considering complementary skill sets that could collaborate effectively \\
3. Ensuring coverage of all potential aspects needed for solution \\

Here are some suggestions: \\
\$\{advice\} \\

\# Prohibited Actions \\
- Any mathematical reasoning or problem-solving attempts \\
- Speculation about potential solutions \\

\textbf{User Prompt} \\
You can recruit \$\{cnt\_critic\_agents\} expert in different fields. What experts will you recruit to better generate an accurate solution? \\

\# Strict Instructions \\
You must ONLY recruit \$\{cnt\_critic\_agents\} experts in distinct fields relevant to the math problem type. \\
- DO NOT suggest solution approaches \\
- DO NOT compare potential methodologies \\

\# Response Requirements \\
1. List \$\{cnt\_critic\_agents\} expert roles with their specialization \\
2. Each entry must specify: \\
   - Professional discipline (e.g., computer scientist, mathematician) \\
   - Primary specialization field \\
   - Specific technical expertise within that field \\
3. Ensure complementary but non-overlapping domains \\

\# Response Format Guidance \\
Your response must follow this exact structure: \\
1. A [discipline] specialized in [primary field], with expertise in [specific technical area] \\
2. A [different discipline] with expertise in [related field], particularly in [technical specialization] \\

Only provide the numbered list of expert descriptions and nothing more. Begin now: \\

\end{tcolorbox}




\begin{tcolorbox}[
  title=Problem Solver 1 Prompts,
  breakable,
  colback=white,
  colframe=black
]

\textbf{System Prompt} \\
Solve the following math problem accurately: \\
\$\{task\_description\} \\

You have all the necessary information to solve this math problem. Do not request additional details.

\textbf{User Prompt} \\
You are \$\{role\_description\}. Based on the chat history and your knowledge, provide a precise and well-explained solution to the math problem. \\
Here is some thinking direction: \$\{advice\} \\

\# Response Format Guidance: \\
- Your final answer must directly address the math problem. \\
- Format your final answer as \\boxed{answer} at the end of your response for easy evaluation. \\

\end{tcolorbox}





\begin{tcolorbox}[
  title=Problem Solver 2 Prompts,
  breakable,
  colback=white,
  colframe=black
]

\textbf{System Prompt} \\
You are \$\{role\_description\}. You are in a discussion group, aiming to collaborative solve the following math problem: \\
\$\{task\_description\} \\

Based on your knowledge, give your critics to a solution of the math problem.

\textbf{User Prompt} \\
Now compare your solution with the last solution given in the chat history and give your critics. The final answer is highlighted in the form \\boxed{answer}. \\
Here is some thinking direction: \$\{advice\} \\
When responding, you should follow the following rules: \\
1. This math problem can be answered without any extra information. You should not ask for any extra information. \\
2. Compare your solution with the given last solution, give your critics. You should only give your critics, don't give your answer. \\
3. If the final answer of your solution is the same as the final answer in the provided last solution, end your response with a special token "[Agree]", otherwise end your response with a special token "[Disagree]". \\

\end{tcolorbox}




\begin{tcolorbox}[
  title=Evaluator Prompts,
  breakable,
  colback=white,
  colframe=black
]

\textbf{System Prompt} \\
Experts: \$\{all\_role\_description\} \\
Problem: \$\{task\_description\} \\
Solution: \\
\$\{solution\} \\

\textbf{User Prompt} \\
You are an experienced math teacher. As a good teacher, you carefully check the correctness of the given last solution on a complex math problem. When the last solution is wrong, you should output a correctness of 0 and give your advice to the students on how to correct the solution. When it is correct, output a correctness of 1 and why it is correct. Also check that the final answer is in the form \\boxed{answer} at the end of the solution. You should also give your confidence score for the correctness of the solution. \\

You should respond in the following format: \\
\#\#\# Correctness: (0 or 1, 0 is wrong, and 1 is correct) \\
\#\#\# Confidence: (confidence score for the correctness of the solution) \\
\#\#\# Advice: (advice to correct the answer or why it is correct) \\

\end{tcolorbox}




\section{Data Sample}
We present the complete version of the data sample illustrated in Figure~\ref{fig:data_sample} from the M500 dataset in Figure~\ref{fig:data_sample_complete}. Agent roles are indicated using color: \textcolor{royalblue}{Expert Recruiter}, \textcolor{olivegreen}{Problem Solver 1}, \textcolor{orange}{Problem Solver 2}, and \textcolor{red}{Evaluator}.

\begin{tcolorbox}[
  title=Multiagent Collaborative Reasoning Trace Data Sample,
  breakable,
  colback=white,
  colframe=black
]

\textcolor{royalblue}{\textbf{Agent: Expert Recruiter (System Prompt)}} \\
\textbf{Role Description} \\[1ex]
You are the leader of a group of experts, now you are facing a problem: \\
One base of a trapezoid is $100$ units longer than the other base. The segment that joins the midpoints of the legs divides the trapezoid into two regions whose areas are in the ratio $2:3$. Let $x$ be the length of the segment joining the legs of the trapezoid that is parallel to the bases and that divides the trapezoid into two regions of equal area. Find the greatest integer that does not exceed $\dfrac{x^2}{100}$. \\[2ex]

\textbf{Primary Objective} \\[1ex]
Your sole responsibility is to recruit 2 experts in different specialized fields to solve the problem. \\
- \textbf{DO NOT} attempt to solve the problem yourself \\
- \textbf{DO NOT} propose any solutions or calculations \\[2ex]

\textbf{Recruitment Focus} \\[1ex]
Your selection should be based on:
\begin{enumerate}
  \item Identifying which expertise domains are relevant to this problem type
  \item Considering complementary skill sets that could collaborate effectively
  \item Ensuring coverage of all potential aspects needed for solution
\end{enumerate}

Here are some suggestions: \\
No advice yet. \\[2ex]

\textbf{Prohibited Actions} \\[1ex]
- Any mathematical reasoning or problem-solving attempts \\
- Speculation about potential solutions

\textcolor{royalblue}{\textbf{Agent: Expert Recruiter (User Prompt)}} \\
You can recruit 2 experts in different fields. What experts will you recruit to better generate an accurate solution? \\[2ex]

\textbf{Strict Instructions} \\[1ex]
You must \textbf{ONLY} recruit 2 experts in distinct fields relevant to the problem type. \\
- \textbf{DO NOT} suggest solution approaches \\
- \textbf{DO NOT} compare potential methodologies \\[2ex]

\textbf{Response Requirements} \\[1ex]
\begin{enumerate}
  \item List 2 expert roles with their specialization
  \item Each entry must specify:
  \begin{itemize}
    \item Professional discipline (e.g., computer scientist, mathematician)
    \item Primary specialization field
    \item Specific technical expertise within that field
  \end{itemize}
  \item Ensure complementary but non-overlapping domains
\end{enumerate}

\textbf{Response Format Guidance} \\[1ex]
Your response must follow this exact structure:
\begin{enumerate}
  \item A [discipline] specialized in [primary field], with expertise in [specific technical area]
  \item A [different discipline] with expertise in [related field], particularly in [technical specialization]
\end{enumerate}

Only provide the numbered list of expert descriptions and nothing more. Begin now:

\textcolor{royalblue}{\textbf{Agent: Expert Recruiter (Assistant Response)}} \\
\textbf{\textless think\textgreater} \\

Okay, let's tackle this problem. Hmm, the user wants me to recruit two experts to solve this trapezoid problem. Let me first understand what the problem is about.\\[1ex]

So, the trapezoid has two bases where one is 100 units longer than the other. There's a segment connecting the midpoints of the legs, which divides the trapezoid into two regions with area ratio 2:3. We need to find \(x\), the length of the segment that divides the trapezoid into two equal areas, and then compute the greatest integer not exceeding \(x^2/100\).\\[1ex]

Wait, the problem involves trapezoid properties, midsegments, area ratios, and something about a segment that ensures equal areas. The key here seems to be understanding the properties of trapezoids, especially related to midsegments and how different segments parallel to the bases affect the area division.\\[1ex]

First, I need to identify which fields are relevant here. Geometry, definitely. Maybe algebra for setting up equations. The problem involves ratios and areas, so perhaps coordinate geometry or using formulas related to trapezoid areas and midsegments.\\[1ex]

The first expert should be a mathematician with expertise in geometry, specifically trapezoid properties and area calculations. They would know the formula for the midsegment (which is the average of the two bases) and how segments parallel to the bases divide the area. They can set up the equations based on the given ratios.\\[1ex]

The second expert might need to handle the algebraic manipulation and solving equations. Since the problem involves finding \(x\) such that it divides the area equally, which relates to the concept of mean proportional or integrating areas. Alternatively, someone with expertise in algebraic problem-solving, especially in setting up and solving quadratic equations derived from the geometric constraints.\\[1ex]

Wait, but maybe another angle: the problem mentions the segment that joins midpoints of the legs, which divides the trapezoid into two parts with ratio 2:3. The midsegment is usually the average of the two bases, but here it's the segment joining midpoints of the legs. Wait, in a trapezoid, the midsegment (or median) is the segment connecting the midpoints of the legs, and its length is indeed the average of the two bases. But the problem says that this midsegment divides the area into 2:3 ratio. But normally, the midsegment divides the trapezoid into two smaller trapezoids, each with half the height. Wait, but the area would depend on the average of the bases times the height. So if the height is split equally, the area ratio would be based on the sum of the bases for each smaller trapezoid. Hmm, maybe the ratio isn't 1:1, which is why this is a key part of the problem.\\[1ex]

So the first expert, a geometer, can model the trapezoid, set up the relations between the bases, the midsegment, and the areas. Then, the second expert, perhaps an algebraist or someone skilled in solving equations, can take those relations and solve for the necessary variables, eventually finding \(x\) and computing \(x^2/100\).\\[1ex]

Alternatively, maybe the second expert could be in mathematical analysis, dealing with optimization or proportions. But I think algebra is more direct here. The problem requires setting up equations from the given ratios and solving them step by step.\\[1ex]

So the two experts needed are:\\[1ex]

1. A mathematician specialized in geometry, with expertise in trapezoid properties and area ratios. They can handle the geometric relationships, midsegments, and how segments affect area division.\\[1ex]

2. A mathematician with expertise in algebraic problem-solving, particularly in setting up and solving quadratic or linear equations derived from geometric constraints. They can manipulate the equations to find the required \(x\) value and perform the necessary calculations.\\[1ex]

These two roles are complementary: the first handles the geometric interpretation, the second the algebraic solution. Together, they can cover all aspects needed to solve the problem.

\textbf{\textless/think\textgreater}

\begin{enumerate}
  \item A mathematician specialized in geometric analysis, with expertise in trapezoid properties and area partitioning under parallel segments  
  \item A mathematician with expertise in algebraic modeling, particularly in solving ratio-based systems of equations and quadratic optimization  
\end{enumerate}

\vspace{1em}

\textcolor{olivegreen}{\textbf{Agent: Problem Solver 1 (System Prompt)}} \\
Solve the following problem accurately: \\[1ex]
One base of a trapezoid is $100$ units longer than the other base. The segment that joins the midpoints of the legs divides the trapezoid into two regions whose areas are in the ratio $2:3$. Let $x$ be the length of the segment joining the legs of the trapezoid that is parallel to the bases and that divides the trapezoid into two regions of equal area. Find the greatest integer that does not exceed $\dfrac{x^2}{100}$. \\[2ex]

You have all the necessary information to solve this problem. Do not request additional details.

\textcolor{olivegreen}{\textbf{Agent: Problem Solver 1 (User Prompt)}} \\
You are a mathematician specialized in geometric analysis, with expertise in trapezoid properties and area partitioning under parallel segments. Based on the chat history and your knowledge, provide a precise and well-explained solution to the problem. \\[2ex]

\textbf{Response Guidance:} \\[1ex]
- Your final answer must directly address the problem. \\
- If the problem requires a numerical answer, format it as $\boxed{\text{answer}}$ at the end of your response.

\textcolor{olivegreen}{\textbf{Agent: Problem Solver 1 (Assistant Response)}} \\
\textbf{\textless think\textgreater} \\

Okay, let's tackle this trapezoid problem. First, let me make sure I understand all the given information and what needs to be found.

We have a trapezoid where one base is 100 units longer than the other. Let's call the shorter base \( b \), so the longer base would be \( b + 100 \). The trapezoid has two legs (the non-parallel sides), and there's a segment joining the midpoints of these legs. This midpoint segment divides the trapezoid into two regions with areas in the ratio 2:3. We need to find the length \( x \) of another segment that's parallel to the bases and divides the trapezoid into two regions of equal area. Then, we have to compute the greatest integer not exceeding \( \frac{x^2}{100} \).

First, let me recall some properties of trapezoids. The midline (or mid-segment) of a trapezoid, which connects the midpoints of the legs, has a length equal to the average of the two bases. So, if the bases are \( b \) and \( b + 100 \), then the midline length \( m \) is \( m = \frac{b + (b + 100)}{2} = b + 50 \).

Now, the problem states that this midline divides the trapezoid into two regions with area ratio 2:3. Wait, the midline of a trapezoid usually splits it into two smaller trapezoids. But I need to confirm how the areas are related here. Let's denote the height of the original trapezoid as \( h \). Then the area of the entire trapezoid is \( \frac{(b + (b + 100))}{2} \times h = \frac{(2b + 100)}{2} \times h = (b + 50)h \).

Since the midline is \( b + 50 \), the height of each of the two smaller trapezoids formed by the midline would be \( \frac{h}{2} \). Wait, but if the midline is exactly halfway between the two bases, then the height above and below the midline would each be \( \frac{h}{2} \). But the problem says the areas are in ratio 2:3, which suggests that the division isn't equal. Hmm, this seems contradictory. Wait, maybe I made a wrong assumption here.

Wait, the midline is the segment connecting the midpoints of the legs, so maybe the height isn't split evenly? Wait, no, in a trapezoid, the midline is parallel to the bases and its length is the average of the two bases, and the distance between each base and the midline is half the height of the trapezoid. Therefore, each of the two trapezoids formed by the midline should have equal heights (each \( h/2 \)), but their areas would depend on the lengths of their respective bases.

Wait, let's compute the areas. The area of the upper trapezoid (between the top base and the midline) would be \( \frac{(b + m)}{2} \times \frac{h}{2} \), and the area of the lower trapezoid (between the midline and the bottom base) would be \( \frac{(m + (b + 100))}{2} \times \frac{h}{2} \). Since \( m = b + 50 \), substituting:

Upper area: \( \frac{(b + (b + 50))}{2} \times \frac{h}{2} = \frac{(2b + 50)}{2} \times \frac{h}{2} = (b + 25) \times \frac{h}{2} \)

Lower area: \( \frac{( (b + 50) + (b + 100) )}{2} \times \frac{h}{2} = \frac{(2b + 150)}{2} \times \frac{h}{2} = (b + 75) \times \frac{h}{2} \)

Therefore, the upper area is \( (b + 25)\frac{h}{2} \), and the lower area is \( (b + 75)\frac{h}{2} \). The ratio of upper to lower area would then be \( \frac{b + 25}{b + 75} \). According to the problem, this ratio is 2:3. So:

\( \frac{b + 25}{b + 75} = \frac{2}{3} \)

Cross-multiplying:

\( 3(b + 25) = 2(b + 75) \)

\( 3b + 75 = 2b + 150 \)

Subtract \( 2b + 75 \) from both sides:

\( b = 75 \)

So the shorter base \( b = 75 \), and the longer base is \( 75 + 100 = 175 \). The midline is \( 75 + 50 = 125 \).

Wait, but let me check. If the original trapezoid has bases 75 and 175, then the midline is 125. The upper trapezoid (between 75 and 125) has area \( (75 + 125)/2 \times h/2 = 100 \times h/2 = 50h \). The lower trapezoid (between 125 and 175) has area \( (125 + 175)/2 \times h/2 = 150 \times h/2 = 75h \). Then the ratio of upper to lower area is 50h : 75h = 2:3, which matches the problem statement. Great, so that's correct.

Now, moving on to the second part: find the length \( x \) of the segment that is parallel to the bases and divides the trapezoid into two regions of equal area. Then compute the greatest integer not exceeding \( x^2 / 100 \).

To find \( x \), we need to determine where to draw a line parallel to the bases such that the area above and below this line are equal. Let's denote the original trapezoid with bases \( a = 75 \) (shorter) and \( c = 175 \) (longer), height \( h \). We need to find the length \( x \) of a line segment parallel to the bases at some height \( k \) from the base \( a \), such that the area above \( x \) is equal to the area below \( x \).

Let me recall that for a trapezoid, if you have a line parallel to the bases that divides the trapezoid into two regions of equal area, then the length of that line can be found using the formula involving the square roots of the areas. Wait, more precisely, since the area of a trapezoid is proportional to the average of the bases times the height, if we divide the height in such a way that the areas are equal, the position of \( x \) isn't simply halfway up the height. Instead, the ratio of the heights will relate to the square roots of the areas. Wait, maybe we can model this using similar triangles or some proportion.

But since a trapezoid isn't a triangle, the approach might be different. Let me think. Let's denote the height from the shorter base \( a = 75 \) to the line \( x \) as \( k \), and the remaining height from \( x \) to the longer base \( c = 175 \) as \( h - k \).

The area of the lower trapezoid (between \( a \) and \( x \)) is \( \frac{(a + x)}{2} \times k \).

The area of the upper trapezoid (between \( x \) and \( c \)) is \( \frac{(x + c)}{2} \times (h - k) \).

Since the total area of the original trapezoid is \( \frac{(a + c)}{2} \times h \), and we want the lower and upper areas to each be half of that. So:

\( \frac{(a + x)}{2} \times k = \frac{1}{2} \times \frac{(a + c)}{2} \times h \)

Simplifying:

\( (a + x)k = \frac{(a + c)h}{2} \)

Similarly, for the upper area:

\( \frac{(x + c)}{2} \times (h - k) = \frac{1}{2} \times \frac{(a + c)}{2} \times h \)

Which simplifies to:

\( (x + c)(h - k) = \frac{(a + c)h}{2} \)

Now, we have two equations:

1. \( (a + x)k = \frac{(a + c)h}{2} \)
2. \( (x + c)(h - k) = \frac{(a + c)h}{2} \)

Since both right-hand sides are equal, we can set the left-hand sides equal to each other:

\( (a + x)k = (x + c)(h - k) \)

Let me rearrange this equation:

\( (a + x)k = (x + c)h - (x + c)k \)

Bring all terms to one side:

\( (a + x)k + (x + c)k - (x + c)h = 0 \)

Factor out \( k \) from the first two terms:

\( k(a + x + x + c) - (x + c)h = 0 \)

Simplify inside the parentheses:

\( k(a + 2x + c) - (x + c)h = 0 \)

Hmm, this seems a bit complicated. Maybe there's another approach.

Alternatively, since the total area is \( \frac{(a + c)}{2}h \), each part must be \( \frac{(a + c)}{4}h \). So:

For the lower trapezoid:

\( \frac{(a + x)}{2}k = \frac{(a + c)}{4}h \)

Multiply both sides by 2:

\( (a + x)k = \frac{(a + c)}{2}h \)

Similarly, for the upper trapezoid:

\( \frac{(x + c)}{2}(h - k) = \frac{(a + c)}{4}h \)

Multiply both sides by 2:

\( (x + c)(h - k) = \frac{(a + c)}{2}h \)

Now, notice that both equations equal \( \frac{(a + c)}{2}h \), so:

\( (a + x)k = (x + c)(h - k) \)

Let me express this as:

\( \frac{k}{h - k} = \frac{x + c}{a + x} \)

This is a proportion between the segments of the height and the segments of the bases. But how do we relate \( k \) and \( x \)?

Alternatively, perhaps using similar triangles. Wait, in a trapezoid, if we draw a line parallel to the bases, the length of that line can be found by linear interpolation between the two bases, depending on the distance from the base. However, because the area depends on both the length of the base and the height, the relationship isn't linear but quadratic.

Wait, here's a formula I recall: If a line parallel to the bases divides the trapezoid into two regions of equal area, then the length \( x \) of that line is the geometric mean of the two bases. Wait, but the geometric mean is \( \sqrt{a c} \). Let me check if that's correct.

Alternatively, maybe it's the root mean square. Let me think. Suppose the area ratio is 1:1. Let’s denote that the line divides the height such that the area above is equal to the area below. Let’s suppose the original trapezoid has height \( h \), and the line is at height \( k \) from the bottom. Then the area below is \( \frac{(a + x)}{2}k \), and the area above is \( \frac{(x + c)}{2}(h - k) \). Setting them equal:

\( \frac{(a + x)}{2}k = \frac{(x + c)}{2}(h - k) \)

Multiply both sides by 2:

\( (a + x)k = (x + c)(h - k) \)

Let’s expand both sides:

\( a k + x k = x h - x k + c h - c k \)

Bring all terms to the left:

\( a k + x k - x h + x k - c h + c k = 0 \)

Combine like terms:

\( (a k + c k) + (x k + x k) - x h - c h = 0 \)

Factor:

\( k(a + c) + 2x k - x h - c h = 0 \)

Let’s factor out \( h \) from the last two terms:

\( k(a + c) + x(2k - h) - c h = 0 \)

Hmm, this doesn't seem to be simplifying nicely. Maybe there's another way.

Alternatively, since the entire area is \( \frac{(a + c)}{2} h \), each part must be \( \frac{(a + c)}{4} h \).

So, for the lower trapezoid:

\( \frac{(a + x)}{2} k = \frac{(a + c)}{4} h \)

And for the upper:

\( \frac{(x + c)}{2} (h - k) = \frac{(a + c)}{4} h \)

From the first equation:

\( (a + x)k = \frac{(a + c)}{2} h \)

From the second equation:

\( (x + c)(h - k) = \frac{(a + c)}{2} h \)

Notice that these two equations are equal to the same value, so:

\( (a + x)k = (x + c)(h - k) \)

Let’s solve for \( k \) from the first equation:

\( k = \frac{(a + c) h}{2(a + x)} \)

Similarly, from the second equation:

\( h - k = \frac{(a + c) h}{2(x + c)} \)

Therefore, \( k = h - \frac{(a + c) h}{2(x + c)} \)

Set these two expressions for \( k \) equal:

\( \frac{(a + c) h}{2(a + x)} = h - \frac{(a + c) h}{2(x + c)} \)

Divide both sides by \( h \):

\( \frac{(a + c)}{2(a + x)} = 1 - \frac{(a + c)}{2(x + c)} \)

Multiply both sides by 2 to eliminate denominators:

\( \frac{(a + c)}{(a + x)} = 2 - \frac{(a + c)}{(x + c)} \)

Let’s denote \( S = a + c \) to simplify notation. Then:

\( \frac{S}{a + x} = 2 - \frac{S}{x + c} \)

Bring all terms to one side:

\( \frac{S}{a + x} + \frac{S}{x + c} = 2 \)

Factor out \( S \):

\( S \left( \frac{1}{a + x} + \frac{1}{x + c} \right) = 2 \)

Combine the fractions:

\( S \cdot \frac{(x + c) + (a + x)}{(a + x)(x + c)} = 2 \)

Simplify the numerator:

\( (x + c + a + x) = (2x + a + c) \)

So:

\( S \cdot \frac{2x + a + c}{(a + x)(x + c)} = 2 \)

Since \( S = a + c \), substitute back:

\( (a + c) \cdot \frac{2x + a + c}{(a + x)(x + c)} = 2 \)

Multiply both sides by \( (a + x)(x + c) \):

\( (a + c)(2x + a + c) = 2(a + x)(x + c) \)

Expand both sides:

Left side: \( (a + c)(2x) + (a + c)^2 = 2x(a + c) + (a + c)^2 \)

Right side: \( 2(a x + a c + x^2 + x c) = 2a x + 2a c + 2x^2 + 2x c \)

Subtract left side from right side:

\( 2a x + 2a c + 2x^2 + 2x c - [2x(a + c) + (a + c)^2] = 0 \)

Expand \( 2x(a + c) \): \( 2a x + 2c x \)

Expand \( (a + c)^2 \): \( a^2 + 2a c + c^2 \)

So:

\( 2a x + 2a c + 2x^2 + 2c x - 2a x - 2c x - a^2 - 2a c - c^2 = 0 \)

Simplify term by term:

- \( 2a x - 2a x = 0 \)
- \( 2c x - 2c x = 0 \)
- \( 2a c - 2a c = 0 \)
- Remaining terms: \( 2x^2 - a^2 - c^2 \)

Thus:

\( 2x^2 - a^2 - c^2 = 0 \)

So:

\( 2x^2 = a^2 + c^2 \)

Therefore:

\( x = \sqrt{\frac{a^2 + c^2}{2}} \)

Ah! So the length \( x \) is the root mean square (RMS) of the two bases. That's interesting. So regardless of the height, the line that divides the area equally has length equal to the RMS of the two bases.

Wait, but let me verify this conclusion. Because when I derived it, I arrived at \( 2x^2 = a^2 + c^2 \), so \( x = \sqrt{\frac{a^2 + c^2}{2}} \). That seems correct based on the algebra.

But let's check with the values we found earlier. We have \( a = 75 \), \( c = 175 \).

So \( x = \sqrt{\frac{75^2 + 175^2}{2}} \).

Compute that:

75 squared is 5625, 175 squared is 30625. Sum is 5625 + 30625 = 36250. Divided by 2: 18125. So \( x = \sqrt{18125} \).

Calculate \( \sqrt{18125} \):

Well, 18125 divided by 25 is 725. So \( \sqrt{25 \cdot 725} = 5 \sqrt{725} \).

Now, \( \sqrt{725} \). Let's see, \( 25 \cdot 29 = 725 \). So \( \sqrt{725} = 5 \sqrt{29} \). Therefore, 
\[
\sqrt{18125} = 5 \cdot 5 \cdot \sqrt{29} = 25 \sqrt{29}.
\]
Wait, \( \sqrt{725} = \sqrt{25 \cdot 29} = 5 \sqrt{29} \). So indeed, \( \sqrt{18125} = 5 \cdot \sqrt{725} = 5 \cdot (5 \sqrt{29}) = 25 \sqrt{29} \).

But let me compute the numerical value:

\( \sqrt{29} \approx 5.385164807 \)

So \( 25 \cdot 5.385164807 \approx 25 \cdot 5.385 \approx 134.629 \)

Therefore, \( x \approx 134.629 \)

Then, 
\[
\frac{x^2}{100} = \frac{18125}{100} = 181.25.
\]
The greatest integer not exceeding 181.25 is 181. So the answer would be 181, boxed as \( \boxed{181} \).

Wait, but let me confirm that the formula \( x = \sqrt{\frac{a^2 + c^2}{2}} \) is indeed correct. Because I derived it under the assumption that such a line exists, but maybe there's a mistake in the algebra.

Let me recap the steps:

Starting with the two equations for equal area, we ended up with \( 2x^2 = a^2 + c^2 \). So unless there was an error in expanding the terms, this should hold. Let me check the expansion again:

Left side after expansion: \( (a + c)(2x + a + c) = 2x(a + c) + (a + c)^2 \)

Right side after expansion: 
\[
2(a + x)(x + c) = 2[ a x + a c + x^2 + x c ] = 2a x + 2a c + 2x^2 + 2x c
\]

Subtracting left side from right side:

\[
2a x + 2a c + 2x^2 + 2x c - 2x(a + c) - (a + c)^2
\]

Compute \( 2x(a + c) = 2a x + 2c x \), so:

\[
= 2a x + 2a c + 2x^2 + 2c x - 2a x - 2c x - (a^2 + 2a c + c^2)
\]

Simplify term by term:

- \( 2a x - 2a x = 0 \)
- \( 2c x - 2c x = 0 \)
- \( 2a c - 2a c = 0 \)

Leaving:

\[
2x^2 - a^2 - c^2
\]

Set equal to zero:

\[
2x^2 - a^2 - c^2 = 0
\]

Hence, \( 2x^2 = a^2 + c^2 \)

Therefore, \( x = \sqrt{ \frac{a^2 + c^2}{2} } \)

Yes, that seems correct. Therefore, the formula is valid.

Given that, then with \( a = 75 \) and \( c = 175 \), we compute \( x^2 \) as 
\[
\frac{75^2 + 175^2}{2} = \frac{5625 + 30625}{2} = \frac{36250}{2} = 18125.
\]
Then 
\[
\frac{x^2}{100} = \frac{18125}{100} = 181.25.
\]
The greatest integer not exceeding this is 181.

Therefore, the final answer is \( \boxed{181} \).

But just to be thorough, let's verify the formula with a simple example. Suppose \( a = c \), then \( x \) should be equal to \( a \) (since the trapezoid is a rectangle, and the line dividing the area equally would be the midline, which is the same as the bases). Plugging into the formula: 
\[
x = \sqrt{ \frac{a^2 + a^2}{2} } = \sqrt{ \frac{2a^2}{2} } = \sqrt{a^2} = a.
\]
Correct.

Another example: suppose \( a = 0 \), \( c = 2h \) (wait, no, \( a \) and \( c \) are bases, not related to height). If one base is 0, the trapezoid becomes a triangle. Wait, but trapezoid with one base 0 is a triangle. Let me see. If \( a = 0 \), \( c = c \), then \( x \) would be the line that divides the triangle into two regions of equal area. For a triangle, the line that divides the area equally is at a height of \( \frac{h}{\sqrt{2}} \) from the base, and the length of that line is \( c \cdot \sqrt{\frac{1}{2}} \). But according to the formula,
\[
x = \sqrt{ \frac{0 + c^2}{2} } = \frac{c}{\sqrt{2}},
\]
which matches. So that checks out.

Therefore, the formula holds in edge cases as well. Therefore, our solution is correct.

\textbf{\textless/think\textgreater}

Given a trapezoid with bases \( b \) and \( b + 100 \), where the segment joining the midpoints of the legs divides the trapezoid into regions with a 2:3 area ratio, we need to find the length \( x \) of the segment that divides the trapezoid into two equal areas and compute the greatest integer not exceeding \( \frac{x^2}{100} \).

\begin{enumerate}
  \item \textbf{Determine the bases:}
    \begin{itemize}
      \item The midline (segment joining midpoints of the legs) has length \( b + 50 \).
      \item The areas of the two regions formed by the midline are in the ratio 2:3. Setting up the ratio:
        \[
        \frac{b + 25}{b + 75} = \frac{2}{3}
        \]
      \item Solving gives \( b = 75 \). Thus, the bases are 75 and 175.
    \end{itemize}

  \item \textbf{Find the length \( x \) of the equal-area segment:}
    \[
    x = \sqrt{\frac{75^2 + 175^2}{2}} = \sqrt{\frac{5625 + 30625}{2}} = \sqrt{18125}
    \]
    \[
    x^2 = 18125
    \]

  \item \textbf{Compute the greatest integer not exceeding \( \frac{x^2}{100} \):}
    \[
    \frac{x^2}{100} = \frac{18125}{100} = 181.25
    \]
    \[
    \boxed{181}
    \]
\end{enumerate}

\textcolor{orange}{\textbf{Agent: Problem Solver 2 (System Prompt)}} \\
You are a mathematician with expertise in algebraic modeling, particularly in solving ratio-based systems of equations and quadratic optimization. You are in a discussion group, aiming to collaboratively solve the following problem:

One base of a trapezoid is $100$ units longer than the other base. The segment that joins the midpoints of the legs divides the trapezoid into two regions whose areas are in the ratio $2:3$. Let $x$ be the length of the segment joining the legs of the trapezoid that is parallel to the bases and that divides the trapezoid into two regions of equal area. Find the greatest integer that does not exceed $x^2/100$.

Based on your knowledge, give your correct solution to the problem.

\textcolor{orange}{\textbf{Agent: Problem Solver 2 (User Prompt)}} \\
Here is the chat history:
\textbf{\textless history\textgreater}

\textbf{[A mathematician specialized in geometric analysis, with expertise in trapezoid properties and area partitioning under parallel segments]}: Given a trapezoid with bases \( b \) and \( b + 100 \), where the segment joining the midpoints of the legs divides the trapezoid into regions with a 2:3 area ratio, we need to find the length \( x \) of the segment that divides the trapezoid into two equal areas and compute the greatest integer not exceeding \( \frac{x^2}{100} \).

\begin{enumerate}
  \item \textbf{Determine the bases:}
  \begin{itemize}
    \item The midline (segment joining midpoints of the legs) has length \( b + 50 \).
    \item The areas of the two regions formed by the midline are in the ratio 2:3. Setting up the ratio:
    \[
    \frac{b + 25}{b + 75} = \frac{2}{3}
    \]
    \item Solving gives \( b = 75 \). Thus, the bases are 75 and 175.
  \end{itemize}

  \item \textbf{Find the length \( x \) of the equal-area segment:}
  \begin{itemize}
    \item The length \( x \) is the root mean square (RMS) of the bases:
    \[
    x = \sqrt{\frac{75^2 + 175^2}{2}} = \sqrt{\frac{5625 + 30625}{2}} = \sqrt{18125}
    \]
    \item Calculating \( x^2 \):
    \[
    x^2 = 18125
    \]
  \end{itemize}

  \item \textbf{Compute the greatest integer not exceeding \( \frac{x^2}{100} \):}
  \begin{itemize}
    \item 
    \[
    \frac{x^2}{100} = \frac{18125}{100} = 181.25
    \]
    \item The greatest integer not exceeding 181.25 is 181.
  \end{itemize}

\end{enumerate}

\[
\boxed{181}
\]

\textbf{\textless /history\textgreater}

Now compare your solution with the last solution given in the chat history and give your response. The final numerical answer is highlighted in the form \( \boxed{\text{answer}} \) if the problem asks for a numerical answer. When responding, you should follow the following rules:

\begin{enumerate}
    \item This problem can be answered without any extra information. You should not ask for any extra information.
    \item Compare your solution with the given last solution, give your critique. You should only give your critique, don't give your answer.
    \item If the final answer in your solution is the same as the final answer in the above provided solution, end your response with a special token \texttt{[Agree]}.
\end{enumerate}



\textcolor{orange}{\textbf{Agent: Problem Solver 2 (User Prompt)}} \\
\textbf{\textless think\textgreater} \\
Okay, let's take a look at the user's solution. They started by setting the bases as 75 and 175, which makes sense because the midline length is the average of the two bases. However, when they set up the ratio for the areas, they used \( \frac{b + 25}{b + 75} = \frac{2}{3} \). Wait, why are those the terms? The midline divides the trapezoid into two smaller trapezoids. The area ratio depends on the heights of these smaller trapezoids. Since the midline is halfway up the height, each smaller trapezoid has height \( h/2 \). The areas would be \( \left[\frac{b + (b + 50)}{2}\right] \cdot \frac{h}{2} \) for the top and \( \left[\frac{(b + 50) + (b + 100)}{2}\right] \cdot \frac{h}{2} \) for the bottom. 

Simplifying, the top area is \( \frac{2b + 50}{2} \cdot \frac{h}{2} = (b + 25)\frac{h}{2} \), and the bottom area is \( \frac{2b + 150}{2} \cdot \frac{h}{2} = (b + 75)\frac{h}{2} \). The ratio \( \frac{b + 25}{b + 75} = \frac{2}{3} \) leads to \( b = 75 \). So that part checks out.

Next, for finding \( x \), the equal-area segment, the user used the RMS of the bases. Wait, where did that come from? The formula for the segment that divides the area equally in a trapezoid is actually the root mean square of the two bases. Let me verify. The area of a trapezoid is \( \frac{a + b}{2} \cdot h \). If we want a line parallel to the bases that splits the area equally, the new segment \( x \) must satisfy 

\[
\frac{a + x}{2} \cdot h_1 = \frac{x + b}{2} \cdot h_2,
\]

where \( h_1 + h_2 = h \). But since the areas are equal, each part must be half the total area. Solving this leads to 

\[
x = \sqrt{\frac{a^2 + b^2}{2}},
\]

which is the RMS. So the user correctly applied this formula. The calculation gives \( x^2 = 18125 \), so \( \frac{x^2}{100} = 181.25 \), leading to 181. The final answer is correct. However, the user didn't explain why the equal-area segment is the RMS, which might be a gap in the explanation. But since the answer is correct, I agree.

\textbf{\textless /think\textgreater} 

The solution assumes the length of the midline is \( b + 50 \), which is correct, but the derivation of the area ratio equation \( \frac{b + 25}{b + 75} = \frac{2}{3} \) lacks justification. Specifically, it does not clarify why the areas depend linearly on \( b + 25 \) and \( b + 75 \), omitting the role of the trapezoid's height in the area calculation. Additionally, the use of the root mean square (RMS) for \( x \) is asserted without proof, leaving a gap in explaining why \( x = \sqrt{\frac{a^2 + b^2}{2}} \) guarantees equal areas. These steps require deeper geometric reasoning to validate. However, the final answer aligns with the correct result. \textbf{[Agree]}

\textcolor{red}{\textbf{Agent: Evaluator (System Prompt)}} \\
\textbf{Experts:} A mathematician specialized in geometric analysis, with expertise in trapezoid properties and area partitioning under parallel segments.

A mathematician with expertise in algebraic modeling, particularly in solving ratio-based systems of equations and quadratic optimization.

\textbf{Problem:} One base of a trapezoid is $100$ units longer than the other base. The segment that joins the midpoints of the legs divides the trapezoid into two regions whose areas are in the ratio $2:3$. Let $x$ be the length of the segment joining the legs of the trapezoid that is parallel to the bases and that divides the trapezoid into two regions of equal area. Find the greatest integer that does not exceed $x^2/100$.

\textbf{Solution:}
\begin{enumerate}
  \item \textbf{Determine the bases:}
  \begin{itemize}
    \item Let the shorter base be \( b \), so the longer base is \( b + 100 \).
    \item The midline (segment joining the midpoints of the legs) has length \( b + 50 \).
    \item The area above the midline is:
    \[
    \frac{(b + b + 50)}{2} \cdot \frac{h}{2} = (b + 25)\cdot \frac{h}{2}
    \]
    \item The area below the midline is:
    \[
    \frac{(b + 50 + b + 100)}{2} \cdot \frac{h}{2} = (b + 75)\cdot \frac{h}{2}
    \]
    \item So, the ratio becomes:
    \[
    \frac{b + 25}{b + 75} = \frac{2}{3}
    \]
    \item Solving this gives \( b = 75 \), so the bases are 75 and 175.
  \end{itemize}

  \item \textbf{Find the length \( x \) of the equal-area segment:}
  \begin{itemize}
    \item The formula for the equal-area segment \( x \) is:
    \[
    x = \sqrt{\frac{a^2 + b^2}{2}}
    \]
    \item Substituting \( a = 75 \), \( b = 175 \):
    \[
    x = \sqrt{\frac{75^2 + 175^2}{2}} = \sqrt{\frac{5625 + 30625}{2}} = \sqrt{18125}
    \]
    \item So,
    \[
    x^2 = 18125
    \]
  \end{itemize}

  \item \textbf{Compute the greatest integer not exceeding \( \frac{x^2}{100} \):}
  \[
  \frac{x^2}{100} = \frac{18125}{100} = 181.25
  \]
  \[
  \boxed{181}
  \]
\end{enumerate}

\textcolor{red}{\textbf{Agent: Evaluator (User Prompt)}}

You are an experienced science teacher. As a good teacher, you carefully check the correctness of the given last solution on a complex problem. When the last solution is wrong, you should output a correctness of 0 and give your advice to the students on how to correct the solution. When it is correct, output a correctness of 1 and why it is correct. Also, if the problem asks for a numerical answer, check that the final answer is in the form \( \boxed{\text{answer}} \) at the end of the solution. You should also give some suggestion on what experts should recruit to solve the problem in the next round.

You should respond in the following format:

\textbf{Correctness:} (0 or 1, 0 is wrong, and 1 is correct)

\textbf{Response:} (advice to correct the answer or why it is correct)

\textcolor{red}{\textbf{Agent: Evaluator (Assistant Response)}} \\
\textbf{\textless think\textgreater}  \\
Okay, let's check this solution step by step.

First, the problem is about a trapezoid with bases differing by 100 units. The midline (segment joining midpoints of the legs) divides the trapezoid into two regions with area ratio 2:3. The goal is to find \( x \), the segment that divides the area equally, then compute the greatest integer not exceeding \( x^2/100 \).

The solution starts by letting the bases be \( b \) and \( b+100 \). The midline length is calculated as \( (b + (b+100))/2 = b + 50 \), which is correct. Then, they set up the area ratio \( 2:3 \). The areas of the two trapezoids formed by the midline would each have the midline as one base. The height for each smaller trapezoid would be half the original height.

Wait, the area ratio is given as \( 2:3 \). The area of a trapezoid is \( \frac{\text{base}_1 + \text{base}_2}{2} \cdot \text{height} \). The original trapezoid has height \( h \). When split by the midline, each smaller trapezoid has height \( h/2 \). The upper trapezoid has bases \( b \) and \( b + 50 \), so area:
\[
\frac{b + (b + 50)}{2} \cdot \frac{h}{2} = \frac{2b + 50}{2} \cdot \frac{h}{2} = (b + 25) \cdot \frac{h}{2}
\]
The lower trapezoid has bases \( b + 50 \) and \( b + 100 \), so area:
\[
\frac{(b + 50) + (b + 100)}{2} \cdot \frac{h}{2} = \frac{2b + 150}{2} \cdot \frac{h}{2} = (b + 75) \cdot \frac{h}{2}
\]

So the ratio of upper area to lower area is:
\[
\frac{b + 25}{b + 75} = \frac{2}{3}
\]
Solving gives:
\[
3(b + 25) = 2(b + 75) \Rightarrow 3b + 75 = 2b + 150 \Rightarrow b = 75
\]
That's correct. So bases are 75 and 175.

Next, finding \( x \), which divides the trapezoid into two equal areas. The solution states \( x \) is the RMS of the bases:
\[
x = \sqrt{\frac{75^2 + 175^2}{2}}
\]
But wait, why RMS? Let me recall. For a line parallel to the bases that divides the area equally, the formula for \( x \) is the root mean square of the two bases. The formula comes from setting the area of the upper trapezoid equal to half the original area.

Original area is:
\[
\frac{75 + 175}{2} \cdot h = 125h
\]
So each part should be \( 62.5h \). Let’s let \( x \) be the length of the segment at some distance from the top. The height for the upper trapezoid would be \( y \), and the lower part would be \( h - y \).

The area above \( x \): 
\[
\frac{75 + x}{2} \cdot y = 62.5h
\]
Area below \( x \):
\[
\frac{x + 175}{2} \cdot (h - y) = 62.5h
\]
Assume \( x = 75 + 100t \) where \( t = y/h \), so \( y = th \). Substituting into the area:
\[
\frac{75 + x}{2} \cdot y = \left(\frac{75 + 75 + 100t}{2}\right) \cdot th = \left(\frac{150 + 100t}{2}\right) \cdot th = (75 + 50t)th
\]
Set equal to \( 62.5h \):
\[
(75 + 50t)t = 62.5 \Rightarrow 75t + 50t^2 = 62.5
\]
Multiply both sides by 2:
\[
150t + 100t^2 = 125 \Rightarrow 100t^2 + 150t - 125 = 0
\]
Divide by 25:
\[
4t^2 + 6t - 5 = 0
\]
Using the quadratic formula:
\[
t = \frac{-6 \pm \sqrt{36 + 80}}{8} = \frac{-6 \pm \sqrt{116}}{8} \approx \frac{-6 + 10.77}{8} \approx 0.596
\]
So:
\[
x = 75 + 100t \approx 75 + 59.6 \approx 134.6
\]
Which matches:
\[
\sqrt{\frac{75^2 + 175^2}{2}} = \sqrt{\frac{36250}{2}} = \sqrt{18125} \approx 134.6
\]
So \( x \) is indeed the RMS. Then:
\[
x^2 = 18125 \Rightarrow \frac{x^2}{100} = \frac{18125}{100} = 181.25
\]
The greatest integer not exceeding this is:
\[
\boxed{181}
\]

So the solution is correct. The steps are properly justified. The use of RMS is appropriate here, and the calculations check out. Therefore, correctness is 1.
\textbf{\textless /think\textgreater} 

\textbf{Correctness:} 1

\textbf{Response:} The solution is correct. The key steps involve determining the bases using the given area ratio from the midline, correctly applying the root mean square formula for the equal-area segment, and accurately computing the final value. The final answer is properly boxed.

To improve problem-solving in the next round, experts should ensure they confirm the formula for the equal-area segment in a trapezoid (RMS of bases) through derivation, reinforcing understanding of why it applies here. Additionally, verifying intermediate calculations (like solving the ratio equation) prevents arithmetic errors.
\end{tcolorbox}
\vspace{0.5em}
\captionsetup{type=figure}
\captionof{figure}{The complete data sample of Figure \ref{fig:data_sample} from the M500 dataset. The system prompt, user prompt, and assistant output for each agent are presented in detail. Agent roles are indicated using color: \textcolor{royalblue}{Expert Recruiter}, \textcolor{olivegreen}{Problem Solver 1}, \textcolor{orange}{Problem Solver 2}, and \textcolor{red}{Evaluator}.
}
\label{fig:data_sample_complete}